% ==========================================================
% 3. OBJETIVOS (Fase 3, Tema 3.3)
% ==========================================================

\section{Objetivos}

\subsection{Objetivo general}

Contribuir a que los estudiantes, la Escuela de Ingeniería de Sistemas y Computación y
Bienestar Universitario de la UPTC cuenten con información propia para decidir qué
medidas de ergonomía y pausas activas promover, con el fin de reducir las molestias
visuales y osteomusculares asociadas al estudio frente a pantallas.

\subsection{Objetivos específicos y metas}

\begin{center}
\begin{longtable}{c L{6.6cm} L{6.6cm}}
\caption{Objetivos específicos y metas} \\
\toprule
\textbf{N.°} & \textbf{Objetivo específico (el qué)} & \textbf{Meta (la cuantificación)} \\
\midrule
\endfirsthead
\toprule
\textbf{N.°} & \textbf{Objetivo específico (el qué)} & \textbf{Meta (la cuantificación)} \\
\midrule
\endhead
1 & Caracterizar a los estudiantes participantes según sexo, semestre y dispositivo principal de estudio. & Reportar la distribución porcentual de los participantes en cada una de esas tres características. \\
\midrule
2 & Determinar la frecuencia de fatiga visual y de molestias osteomusculares reportadas por los estudiantes. & Reportar la proporción de participantes con fatiga visual y con molestias osteomusculares, y esta última por región corporal. \\
\midrule
3 & Identificar el tiempo diario de exposición a pantallas y los hábitos posturales autorreportados durante las jornadas de estudio. & Reportar el promedio y la dispersión de las horas diarias de exposición, y la frecuencia de cada hábito postural indagado. \\
\midrule
4 & Establecer la relación entre el tiempo de exposición a pantallas y los hábitos posturales, y la frecuencia de fatiga visual y de molestias osteomusculares. & Reportar, para cada par de variables, el coeficiente de asociación y su intervalo de confianza. \\
\bottomrule
\end{longtable}
\end{center}

\noindent El tamaño de la muestra y el instrumento con que se medirá cada variable se
definen en el diseño metodológico (sección 8).
